\documentclass{article}
\usepackage{amsmath}
\usepackage{graphicx}
\usepackage{amssymb}
\usepackage{bm}

\title{Lecture 23: Mode splitting and coupling}
\author{David Al-Attar, Michaelmas term 2024}
\date{}

\begin{document}

\maketitle

\section*{Outline and Motivation}
In this final lecture on seismology we consider the free oscillations of the Earth in further detail. Last time we established the general theory, and saw how eigenfrequencies can be observed from the Fourier transforms of sufficiently long seismograms. Today we first consider the eigenvalue problem in a spherically symmetric earth model for which the eigenfunctions can be calculated with relative ease. An important point is that the eigenspaces in a spherically symmetric model are degenerate, each having dimension $2l+1$ for some integer degree $l\ge0$. This material depends on the properties of vector spherical harmonics, but you are not expected to learn all the details. We then show using first-order perturbation theory how in the real Earth this degeneracy is split by rotation and lateral heterogeneity. To conclude, an overview is given of radial and laterally variations of density within the Earth, with some current research questions highlighted.

\section*{Vector spherical harmonics}
To discuss the eigenfunctions of a spherically symmetric earth model we need to review spherical harmonic expansions of scalar and vector functions. The emphasis here is only on establishing notations that will be useful later. You will not be expected to carry out any detailed calculations with either scalar or vector spherical harmonics.

Using spherical polar co-ordinates $(r, \theta, \varphi)$ the gradient operator is given by
\begin{equation}
\nabla=\hat{\mathbf{r}}\frac{\partial}{\partial r}+\frac{1}{r}\hat{\boldsymbol{\theta}}\frac{\partial}{\partial\theta}+\frac{1}{r\sin\theta}\hat{\boldsymbol{\varphi}}\frac{\partial}{\partial\varphi},
\end{equation}
where, for example, $\hat{\mathbf{r}}$ is a unit vector in the radial direction. This expression can be conveniently written
\begin{equation}
\nabla=\hat{\mathbf{r}}\frac{\partial}{\partial r}+\frac{1}{r}\nabla_{1},
\end{equation}
which acts to define the tangential gradient operator $\nabla_{1}$ for functions defined on the unit-sphere, $\mathbb{S}^{2}$. In terms of $\nabla_{1}$ we can define related operators $\hat{\mathbf{r}}\times\nabla_{1}$ and $\nabla_{1}^{2}$. Up to suitable scale-factors, the first is the quantum mechanical angular momentum operator, $\hat{\mathbf{L}}$, and the second the total squared angular momentum, $\hat{L}^{2}$.

Spherical harmonic functions are simultaneous eigenfunctions of the commuting and self-adjoint operators $\nabla_{1}^{2}$ and $\frac{1}{i}\frac{\partial}{\partial\varphi}$, the latter being proportional to the z-component of $\hat{\mathbf{r}}\times\nabla_{1}$. The defining relations are
\begin{equation}
\nabla_{1}^{2}Y_{lm}=-l(l+1)Y_{lm}; \quad \frac{1}{i}\frac{\partial Y_{lm}}{\partial\varphi}=mY_{lm};
\end{equation}
where $l$ is the degree of the spherical harmonic which takes non-negative integer values, while for fixed $l$ the order, m, is an integer such that $-l\le m\le l$. Note that the eigenspace of $\nabla_{1}^{2}$ with degree l is $(2l+1)$-fold degenerate, this being characteristic of the eigenspaces of rotationally invariant operators. The spherical harmonics satisfy the orthonormality relation
\begin{equation}
\int_{\mathbb{S}^{2}}Y_{lm}^{*}Y_{l'm'}dS=\delta_{ll'}\delta_{mm'}.
\end{equation}
They also form a complete basis for suitably regular function f on $\mathbb{S}^{2}$, and hence we can form the spherical harmonic expansion
\begin{equation}
f=\sum_{l=0}^{\infty}\sum_{m=-l}^{l}f_{lm}Y_{lm},
\end{equation}
where the expansion coefficients are
\begin{equation}
f_{lm}=\int_{\mathbb{S}^{2}}Y_{lm}^{*}f~dS.
\end{equation}
Usually we will write such spherical harmonic expansions simply as $\sum_{lm}f_{lm}Y_{lm}$ leaving implicit the range of the degree and order. To work with vector fields it is convenient to introduce the vector spherical harmonics
\begin{equation}
\mathbf{A}_{lm}=Y_{lm}\hat{\mathbf{r}}, \quad \mathbf{B}_{lm}=\nabla_{1}Y_{lm}, \quad \mathbf{C}_{lm}=-\hat{\mathbf{r}}\times\nabla_{1}Y_{lm}
\end{equation}
It is clear that these vector fields are pairwise orthogonal, while we also have the relations
\begin{align}
\int_{\mathbb{S}^{2}}\mathbf{A}_{l'm'}^{*}\cdot \mathbf{A}_{lm}dS&=\delta_{ll'}\delta_{mm'}, \\
\int_{\mathbb{S}^{2}}\mathbf{B}_{l'm'}^{*}\cdot \mathbf{B}_{lm}dS&=l(l+1)\delta_{ll'}\delta_{mm'}, \\
\int_{\mathbb{S}^{2}}\mathbf{C}_{l'm'}^{*}\cdot \mathbf{C}_{lm}dS&=l(l+1)\delta_{ll'}\delta_{mm'},
\end{align}
where $\cdot$ denotes the usual real inner product $\mathbf{a}\cdot \mathbf{b}=a_{i}b_{i}$ for vectors. Note that $\mathbf{B}_{lm}$ and $\mathbf{C}_{lm}$ could be scaled such that these functions are orthonormal, but this is not necessary. It can be shown that the vector spherical harmonics are complete for vector fields defined on $\mathbb{S}^{2}$, this meaning that such a $\mathbf{u}$ can be expanded as
\begin{equation}
\mathbf{u}=\sum_{lm}(U_{lm}\mathbf{A}_{lm}+V_{lm}\mathbf{B}_{lm}+W_{lm}\mathbf{C}_{lm}) ,
\end{equation}
with coefficients $U_{lm}$, $V_{lm}$ and $W_{lm}$ that can be expressed in terms of $\mathbf{u}$ using the above orthogonality relations.

\subsection*{Spheroidal and toroidal modes}
In a non-rotating earth model, we recall that the eigenvalue problem takes the form
\begin{equation}
-\omega^{2}\langle \mathbf{w}|P|\mathbf{s}\rangle+\langle \mathbf{w}|H|\mathbf{s}\rangle=0
\end{equation}
where $\omega$ is the eigenfrequency, $\mathbf{s}$ the eigenfunction, and $\mathbf{w}$ is an arbitrary test function. When the material parameters $\rho$, and $A_{ijkl}$ are spherically symmetric\footnote{The most general elastic tensor that is consistent with spherical symmetry is transversely isotropic with a radial symmetry axis. There are then five elastic modulii, each of which can vary with radius.}, it can be shown that the possible eigenfunctions take two possible types known as toroidal and spheroidal modes. This can be shown either through a direct but lengthy calculation using spherical harmonics, or using representation theory for the rotation group. For reasons of both time and difficulty, we here just summarise the eventual results.

First there are toroidal modes which have eigenfunctions of the form
\begin{equation}
\mathbf{s}(r,\theta,\varphi)=W(r)\mathbf{C}_{lm}(\theta,\varphi),
\end{equation}
with W a function of radius that depends on $l$ but not m. The corresponding eigenfrequency is also independent of m, and hence such modes lie in $(2l+1)$-fold degenerate eigenspaces. It can be shown that toroidal modes are associated with no perturbation to the gravitational potential, and also that their eigenfunctions are non-zero only in solid parts of the earth model.

The second type of eigenfunction is the spheroidal modes that take the form
\begin{equation}
\mathbf{s}(r,\theta,\varphi)=U(r)\mathbf{A}_{lm}(\theta,\varphi)+V(r)\mathbf{B}_{lm}(\theta,\varphi),
\end{equation}
where U and V are functions of radius that depend on $l$ but not m. Again, such an eigenfunction lies in a $(2l+1)$-fold degenerate eigenspace. Spheroidal modes are associated with gravitational perturbations, and their eigenfunctions can be non-zero in both solid and fluid regions of the earth model.

Let us consider toroidal modes further. If we put eq. (13) into the kinetic energy form, then due to the orthogonality of vector spherical harmonics, we can only get a non-zero contribution from the part of the test function proportional to $\mathbf{C}_{lm}$. The same turns out to be true for the potential energy form, and hence it is sufficient to consider test functions
\begin{equation}
\mathbf{w}(r,\theta,\varphi)=W'(r)\mathbf{C}_{lm}(\theta,\varphi)
\end{equation}
Using eq. (8) it is then easy to obtain
\begin{equation}
\langle \mathbf{w}|P|\mathbf{s}\rangle=\zeta^{2}\int_{0}^{b}\rho W'^{*}Wr^{2}dr,
\end{equation}
where we have set $\zeta=\sqrt{l(l+1)}$ and have assumed the earth model is entirely solid. In a similar manner, it can be shown that
\begin{equation}
\langle \mathbf{w}|H|\mathbf{s}\rangle=\zeta^{2}\int_{0}^{b}\mu(r\frac{dW'^{*}}{dr}-W'^{*})(r\frac{dW}{dr}-W)dr+\zeta^{2}(\zeta^{2}-2)\int_{0}^{b}\mu W'^{*}W~dr.
\end{equation}
Note that derivation of this result is not simple, as you need to know how to differentiate vector spherical harmonics. Nonetheless, the end result is that the eigenvalue problem for a toroidal mode of degree $l$ reduces to the form
\begin{equation}
-\omega^{2}\int_{0}^{b}\rho W'^{*}Wr^{2}dr+\int_{0}^{b}\mu(r\frac{dW'^{*}}{dr}-W'^{*})(r\frac{dW}{dr}-W)dr+(\zeta^{2}-2)\int_{0}^{b}\mu W'^{*}W~dr=0 ,
\end{equation}
where W is the radial eigenfunction, and $W'$ acts as the test function. Through a relatively simple calculation, it can be shown that this is an ODE-eigenvalue problem of Sturm-Lioville type. In the case of a homogeneous planet closed-form solutions can be obtained using spherical Bessel functions, though such a solution is of little practical use. More generally, the resulting ODE eigenvalue problem can be solved numerically, or the weak form directly converted into an algebraic eigenvalue problem using something known as a Galerkin discretisation. For our purposes, the key result is that for each degree $l\ne0$ there exist a countable infinity of toroidal modes, with these different solutions being labelled by an overtone number, $n=0,1,2,...$.

For spheroidal modes a similar approach can be applied, with the eigenvalue problem reducing for each degree, $l$, to one involving the radial eigenfunctions U and V. Solution of this latter problem is considerably more involved than that for toroidal modes, in part due to the need to model the associated gravitational potential perturbations. Nonetheless, numerical solutions of this problem have been routinely determined for decades.

We now summarise the results of this section. Within a spherically symmetric earth model there are two families of normal modes: toroidals and spheroidal. We can introduce an index, $q$, to distinguish these with, say, $q=1$ meaning spheroidal, and $q=2$ toroidal\footnote{Along with further values if we need to distinguish toroidals in different solid parts of the model such as the mantle and inner core.}. Having fixed the mode type, q, we can then consider different degrees, $l$. For each such degree there is an infinite set of eigenfrequencies and radial eigenfunctions labelled by their overtone number, n. Finally, we then have the order, $m$, of a mode, this labelling an orthogonal basis for the $(2l+1)$-fold degenerate eigenspace. In this manner, we see that four integers $(q,l,n,m)$ are sufficient to label all the eigenfunctions of a spherically symmetric earth model, while the corresponding eigenfrequencies depend only on (q,l,n). As a useful shorthand, we can collapse the triplet (q, l, n) into a single index, k, and hence write $\omega_{k}$ for the eigenfrequencies and $|km\rangle$ for the eigenfunctions. Indeed, this is precisely the notation used in the last lecture. When needed, however, we can also write $|qlnm\rangle$ to explicitly characterise an eigenfunction. As a final bit of terminology, for each (q,l,n) we call the eigenspace a multiplet, while individual elements of its basis, as labelled by m, are known as singlets. We might then, for example, talk about a degree two toroidal multiplet which is comprised of five singlets. A common labelling scheme for toroidal multiplets is to write ${}_{n}T_{l}$ for the multiplet of degree $l$ and overtone number n. Similarly, spheroidal multiplets are written ${}_{n}S_{l}$.

\subsection*{Splitting of a multiplet due to rotation and lateral heterogeneity}
Within a spherically symmetric earth model, the normal modes come in $(2l+1)$-fold degenerate multiplets, with $l$ the degree which is a non-negative integer. Such degeneracy is a direct consequence of the spherical symmetry of the problem, and analogous results are seen in other parts of physics. In particular, you know that the energy eigenspaces for the hydrogen atom are similarly degenerate, though note that here there is additional degeneracy associated with the "hidden symmetry" within central force problems due to the Runge-Lenz vector. If such a spherically symmetric system undergoes an aspherical perturbation we expect that this degeneracy will be broken, with the singlets that comprise the multiplet splitting apart. This is precisely what happens in the Zeeman-effect for a hydrogen atom, and we will apply very similar mathematics to examine what rotation and laterally heterogeneities do to the free oscillations of the Earth.

We consider a spherically symmetric earth model along with the non-rotating eigenvalue problem
\begin{equation}
-\omega^{2}\langle \mathbf{w}|P_{0}|\mathbf{s}\rangle+\langle \mathbf{w}|H_{0}|\mathbf{s}\rangle=0.
\end{equation}
where we note that the subscript 0 has been added to the kinetic and potential energy forms. A multiplet of this earth model is identified by the triplet (q,l,n), and we write $\omega_{k}$ and $|km\rangle$ for its eigenfrequency and eigenfunctions using the simplified notations discussed above. We then have, in particular, that
\begin{equation}
-\omega_{k}^{2}\langle \mathbf{w}|P_{0}|km\rangle+\langle \mathbf{w}|H_{0}|km\rangle=0,
\end{equation}
for all test functions $\mathbf{w}$. Suppose now that we add a laterally heterogeneous perturbation to the model such that the kinetic energy form can be written
\begin{equation}
\langle \mathbf{w}|P|\mathbf{s}\rangle=\langle \mathbf{w}|P_{0}|\mathbf{s}\rangle+s\langle \mathbf{w}|P_{1}|\mathbf{s}\rangle+\cdot\cdot\cdot,
\end{equation}
where s is a perturbation parameter. Using a similar notation for the potential energy form, we wish to consider the eigenvalue problem
\begin{equation}
-\omega^{2}[\langle \mathbf{w}|P_{0}|\mathbf{s}\rangle+s\langle \mathbf{w}|P_{1}|\mathbf{s}\rangle+\cdot\cdot\cdot]+[\langle \mathbf{w}|H_{0}|\mathbf{s}\rangle+s\langle \mathbf{w}|H_{1}|\mathbf{s}\rangle+\cdot\cdot\cdot]=0.
\end{equation}
We assume that there is an eigenfrequency for this problem close to $\omega_{k}$, and hence write
\begin{equation}
\omega=\omega_{k}+s\delta\omega+\cdot\cdot\cdot
\end{equation}
while for the eigenfunctions we assume
\begin{equation}
\mathbf{s}=\sum_{m=-l}^{l}a_{m}|km\rangle+s~\delta\mathbf{s}+\cdot\cdot\cdot
\end{equation}
where the $a_{m}$ are coefficients to be determined. Note here that we are allowing for the zeroth-order eigenspace to be degenerate, this method having been applied earlier in the course when considering shear wave splitting. Putting these expansions into the eigenvalue problem, we find at order one that
\begin{equation}
\sum_{m=-l}^{l}[-2\omega_{k}\delta\omega\langle \mathbf{w}|P_{0}|km\rangle-\omega_{k}^{2}\langle \mathbf{w}|P_{1}|km\rangle+\langle \mathbf{w}|H_{1}|km\rangle]a_{m} -\omega_{k}^{2}\langle \mathbf{w}|P_{0}|\delta\mathbf{s}\rangle+\langle \mathbf{w}|H_{0}|\delta\mathbf{s}\rangle=0.
\end{equation}
If we take the test function, $\mathbf{w}$, to equal $|km'\rangle$ then the terms involving $\delta\mathbf{s}$ cancel due to the self-adjointness of the forms and the zeroth-order eigenvalue problem, and hence we are left with
\begin{equation}
\frac{1}{2\omega_{k}}\sum_{m=-l}^{l}[-\omega_{k}^{2}\langle km'|P_{1}|km\rangle+\langle km'|H_{1}|km\rangle]a_{m}=\delta\omega a_{m'}
\end{equation}
where we have used the orthonormality of the singlets. Such an equation is to hold for all $-l\le m'\le l$ and hence we have arrived at a $(2l+1)$-dimensional Hermitian eigenvalue problem for $\delta\omega$ along with the coefficients, $a_{m}$, of the zeroth-order eigenfunction. Assuming that the matrix elements $\langle km'|P_{1}|km\rangle$ and $\langle km'|H_{1}|km\rangle$ can be calculated, such an eigenvalue problem can be readily solved either in closed-form for small $l$, or otherwise on a computer. The important qualitative conclusion of this analysis is that there will generically be $2l+1$ distinct eigenfrequency perturbations $\delta\omega$, this showing that the multiplet is split into the appropriate number of non-degenerate eigenspaces.

The above argument can be readily extended to allow for rotation of the earth model as a first-order perturbation. Indeed, we need only generalise eq. (22) to become
\begin{equation}
-\omega^{2}[\langle \mathbf{w}|P_{0}|\mathbf{s}\rangle+s\langle \mathbf{w}|P_{1}|\mathbf{s}\rangle+\cdot\cdot\cdot]+i\omega[s\langle \mathbf{w}|W|\mathbf{s}\rangle] +[\langle \mathbf{w}|H_{0}|\mathbf{s}\rangle+s\langle \mathbf{w}|H_{1}|\mathbf{s}\rangle+\cdot\cdot\cdot]=0.
\end{equation}
A similar expansion to first-order then leads to the Hermitian eigenvalue problem
\begin{equation}
\frac{1}{2\omega_{k}}\sum_{m=-l}^{l}[-\omega_{k}^{2}\langle km'|P_{1}|km\rangle+i\omega_{k}\langle km'|W|km\rangle+\langle km'|H_{1}|km\rangle]a_{m}=\delta\omega a_{m'}
\end{equation}
where we see on the left hand side an additional set of matrix elements associated with the Coriolis force. Note also that terms associated with the centrifugal force need to be incorporated into the perturbation to the potential energy form.

\subsection*{Matrix elements and selection rules}
The matrix elements in eq. (27) associated with kinetic energy take the explicit form
\begin{equation}
\langle km'|P_{1}|km\rangle=\int_{M}\delta\rho|km'\rangle^{*}\cdot|km\rangle d^{3}x,
\end{equation}
where $\delta\rho$ is the density perturbation\footnote{If aspherical perturbations to the surface or internal boundaries were considered there would also be additional terms, but we neglect these for simplicity.}. Supposing that it is a toroidal mode we are dealing with, then we know that
\begin{equation}
|km\rangle=W_{k}\mathbf{C}_{lm}
\end{equation}
for some radial eigenfunction independent of m. We could also choose to expand the density perturbation in spherical harmonics as
\begin{equation}
\delta\rho=\sum_{st}\delta\rho_{st}Y_{st}
\end{equation}
where we write (s, t) for the degree and order of structural perturbations. Using these results, the matrix element then reduces to
\begin{equation}
\langle km'|P_{1}|km\rangle=\sum_{st}\int_{0}^{b}\left(\int_{\mathbb{S}^{2}}Y_{st}\mathbf{C}_{lm'}^{*}\cdot \mathbf{C}_{lm}dS\right)\delta\rho_{st}W_{k}^{*}W_{k}r^{2}dr.
\end{equation}
Interestingly, we see that singlet orders m and $m'$ enter into the matrix element only through the surface integral
\begin{equation}
\int_{\mathbb{S}^{2}}Y_{st}\mathbf{C}_{lm'}^{*}\cdot \mathbf{C}_{lm}dS,
\end{equation}
which involves three sets of spherical harmonics, two of which have been suitably differentiated. It turns out that such integrals can be evaluated in closed-form. Indeed, up to notation they are the Clebsch-Gordon coefficients occurring in the quantum theory of angular momentum. There is not time to enter deeply into this topic, but we will make a few observations. First, we recall that the spherical harmonic $Y_{lm}$ can takes the form
\begin{equation}
Y_{lm}(\theta,\varphi)=X_{lm}(\theta)e^{im\varphi},
\end{equation}
where $X_{lm}$ can be expressed in terms of associated Legendre polynomials. As a result, the overall azimuthal dependence of the integrand is proportional to $e^{i(-m'+t+m)\varphi}$, and such a term will integrate to zero unless
\begin{equation}
-m'+t+m=0.
\end{equation}
This is a selection rule, and gives a condition for non-trivial interactions between specific singlets within the multiplet under a structural perturbation of order t. Similar ideas apply for other matrix elements, including those formed from distinct multiplets that arise in applications such as quasi-degenerate perturbation theory that is required when a group of multiplets have similar eigenfrequencies. A result known as the Wigner-Eckart theorem shows that the general form of such a matrix element is
\begin{equation}
\langle k'm'|P_{1}|km\rangle=\sum_{st}(-1)^{m}4\pi\nu_{l'}\nu_{s}\nu_{l}\begin{pmatrix} l' & s & l \\ -m' & t & m \end{pmatrix}\langle k'|P_{1}^{st}|k\rangle,
\end{equation}
where we have defined $\nu_{l}=\sqrt{\frac{2l+1}{4\pi}}$ and the term $\langle k'|P_{1}^{st}|k\rangle$ is known as a reduced matrix element which depends on k, k' and the structural perturbations at degree s and order t. The array in this expression is a Wigner 3-j symbol, this being a numerical coefficient occurring in expression such as
\begin{equation}
\int_{\mathbb{S}^{2}}Y_{l'm'}^{*}Y_{st}Y_{lm}dS=(-1)^{m}4\pi\nu_{l'}\nu_{s}\nu_{l}\begin{pmatrix} l' & s & l \\ -m' & t & m \end{pmatrix}\begin{pmatrix} l' & s & l \\ 0 & 0 & 0 \end{pmatrix} .
\end{equation}
The point here is not the details of the formulae, but only that such expression exist and that they contain within them a host of selection rules which govern how the multiplets and singlets interact due to lateral structure. To give a further example, when $k'=k$, then it is found that only even degree structural perturbations contribute to the matrix elements. This applies, in particular, to the splitting calculations of the previous section. To get a sense for why this holds, we recall that a parity transformation takes the form $\mathbf{x}\mapsto-\mathbf{x}$. Under such a transformation, it is clear that the value of the integral in eq.(37) does not change. It can be shown, however, that a spherical harmonic of degree $l$ is multiplied by the factor $(-1)^{l}$ under a parity transformation. Thus we arrive at the equality
\begin{equation}
\int_{\mathbb{S}^{2}}Y_{l'm'}^{*}Y_{st}Y_{lm}dS=(-1)^{l+s+l'}\int_{\mathbb{S}^{2}}Y_{l'm'}^{*}Y_{st}Y_{lm}dS,
\end{equation}
from which we conclude that the integral vanishes unless $l+s+l'$ is even. In situations where $l=l'$ this condition reduces to the structural degree, s, being even. As a final example, it can be shown that multiplets of degree $l$ and $l'$ can be coupled by heterogeneity of degree s if and only if
\begin{equation}
|l-l'|\le s\le l+l'.
\end{equation}
We will not prove this result, but its form should be familiar from the addition of angular momentum in quantum mechanics.

\subsection*{The diagonal sum rule and spherically symmetric density variations}
Returning to the first-order splitting calculations earlier, we were led to a Hermitian eigenvalue problem of the form
\begin{equation}
\frac{1}{2\omega_{k}}\sum_{m=-l}^{l}[-\omega_{k}^{2}\langle km'|P_{1}|km\rangle+i\omega_{k}\langle km'|W|km\rangle+\langle km'|H_{1}|km\rangle]a_{m}=\delta\omega a_{m'}
\end{equation}
The perturbations to the singlet eigenfrequencies are equal to the eigenvalues of the matrix on the left hand side. We will show that this matrix has zero-trace, and hence the singlet eigenfrequency perturbations sum to zero. To see why this holds, we recall the form of a general matrix element in eq.(36). In the case that $k=k'$, the trace of such a matrix element will vanish if the following identity holds
\begin{equation}
\sum_{m=-l}^{l}\begin{pmatrix} l & s & l \\ -m & t & m \end{pmatrix}=0.
\end{equation}
To see that this is true, we use eq. (37) to obtain
\begin{equation}
\sum_{m=-l}^{l}\int_{\mathbb{S}^{2}}Y_{lm}^{*}Y_{st}Y_{lm}dS=(-1)^{m}4\pi\nu_{l'}\nu_{s}\nu_{l}\sum_{m=-l}^{l}\begin{pmatrix} l & s & l \\ -m & t & m \end{pmatrix}\begin{pmatrix} l & s & l \\ 0 & 0 & 0 \end{pmatrix} .
\end{equation}
The left hand side can be simplified using the spherical harmonic addition theorem, a special case of which implies that
\begin{equation}
\sum_{m=-l}^{l}Y_{lm}^{*}(\theta,\varphi)Y_{lm}(\theta,\varphi)=\frac{2l+1}{4\pi}.
\end{equation}
It follows that on the left hand side we have the integral of $Y_{st}$ against a constant function, and so the result vanishes unless $s=t=0$.

This result is known as the diagonal sum rule and was obtained by Freeman Gilbert in 1971. It played an important role in early studies of spherical Earth structure using observations of free oscillation. In the previous lecture we saw that eigenfrequencies can be measured from the amplitude spectra of sufficiently long seismograms. The first such observations were made following the 1960 Great Chilean earthquake which remains the largest event that has occurred during the instrumental period. By this time there existed good models of the Earth's spherically symmetric velocity structure obtained from travel time observations. There also existed plausible models of radial density variations that were built using knowledge of the Earth's mass, moments of inertia, and simple thermodynamic models of density variations in a homogeneous solids compressed under its own weight. As a consequence, the first spectra observed could be immediately compared to eigenfrequencies calculated in the different earth models available.

An example of this can be seen in Fig 3. Here two reasonably clear resonance peaks can be identified, with their centre frequencies measured. Also shown on the plot are the predicted eigenfrequencies from two different spherical earth models for the same multiplet. It is notable, firstly, that these predictions are in reasonable agreement with an observed eigenfrequency. But also that while in a spherical earth model we would expect only one eigenfrequency in this range, there seems to be at least two in the data. The latter fact is due to splitting associated with rotation and lateral heterogeneity. This fact was immediately recognised, and motivated George Backus and Freeman Gilbert to publish in (1961) the first theoretical treatment of the splitting of free oscillations due to rotation. Indeed, as we have seen, such calculations mirror almost exactly those known from quantum mechanics for decades at that point.

During these early days of seismology, the main emphasis was on better constraining the Earth's spherically averaged structure. With travel time observations this could be easily done. For many source-receiver pairs the travel time of different seismic phases could be measured. This information could then be reduced to plots of travel-time against epicentral angle. Due to lateral heterogeneity, along with observational error, there was some scatter in these travel-time curves, but this could be averaged out. Note that the latter step is not rigorously justified, but it was at least easy to do. With free oscillations, however, the situation was quite different. What we observe are split eigenfrequencies of the singlets, but what we really want is the eigenfrequency of the associated multiplet prior to this splitting. Supposing that all the singlet frequencies for a split multiplet have been measured, Gilbert's diagonal sum rule provides an immediate answer. We need only take the arithmetic average of the split eigenfrequencies to obtain the degenerate eigenfrequency corresponding to the Earth's spherically averaged structure. For some low frequency and low degree multiplets it is indeed possible to identify all the singlet eigenfrequencies from a single spectrum. In other cases the singlet frequencies cannot be seen clearly in any one spectrum, but by combining observations from many different earthquakes and stations the singlet frequencies can be recovered using various signal processing methods.

Combining such observations with the diagonal sum rule, a large data base of eigenfrequencies corresponding to the Earth's spherically averaged structure had been determined by the late 1970s. Taken along with travel time information, this allowed existing earth models to be substantially refined, and, in particular, for radial density variations to be found. The resulting density structure is shown in Fig. 4. These observed densities can be compared with predictions made by mineral physicists subject to assumptions about the Earth's interior composition. For example, it has been shown that the density of the outer core is significantly lighter than expected if it were comprised of an iron-nickel alloy as is suggested observations of meteorites. This discrepancy is telling us something about how planets form and differentiation. As yet the identity of this "light element" in the core is uncertain, with oxygen or silicon being perhaps the most likely candidates. But what is worth emphasising is that such discussions are only possible due to the highly accurate constraints on the interior structure provided by seismology.

\subsection*{Lateral density variations and mode coupling}
As noted in Lectures 18 and 19, recent work in seismology is focused on lateral heterogeneity in the mantle, this being done to better understand mantle convection. In Fig. 5 images from a range of S-wave velocity models are shown at different depths within the Earth, and a broad agreement can be seen. In particular, all models show at the base of the mantle two large regions with low shear wave velocity, one below the Pacific, and the other below Africa. On the assumption that these velocity variations reflect temperature variations, then we conclude that these so-called Large Low Shear Velocity Provinces (LLSVPs) are hotter than the surrounding mantle. If this is true, then we would also expect that these regions are lighter than the ambient mantle, and hence will experience an upwards buoyancy force. When LLSVPs were first identified in the early 90s such an interpretation was taken for granted, and they were known as "super plumes". A range of subsequent work, however, has brought into question this conclusion. For example, joint S- and P-wave tomographic models have shown that the LLSVPs seem are associated with a positive perturbation to the bulk sound speed
\begin{equation}
\gamma=\sqrt{\frac{\kappa}{\rho}}.
\end{equation}
Note that this is not a speed of any elastic wave, but rather a parameter that is useful because it depends only on the bulk and not shear modulus. If the lateral variations are caused by temperature, then we would expect that the shear and bulk sound velocities would positively correlate. To see this, we assume that temperature variations are sufficiently small to allow us to write
\begin{equation}
\delta\beta=\frac{\partial\beta}{\partial T}\delta T, \quad \delta\gamma = \frac{\partial\gamma}{\partial T}\delta T.
\end{equation}
where $\delta T$ is the temperature perturbation. The relevant partial derivatives can be estimated from mineral physics, and are negative for all crystalline materials, and hence
\begin{equation}
\frac{\delta\gamma}{\delta\beta}=\frac{\frac{\partial\gamma}{\partial T}}{\frac{\partial\beta}{\partial T}}>0.
\end{equation}
as required. The fact that shear and bulk sound speeds seem to anti-correlate in the lower mantle suggests that temperature is not the only cause for these lateral variations. The other possible explanation is that compositional variations in the mantle play an important role in its dynamics. To see how this might reconcile the observations, we write X for a compositional variables. The velocity perturbations can then be expressed in terms of temperature and composition in the form
\begin{equation}
\delta\beta=\frac{\partial\beta}{\partial T}\delta T+\frac{\partial\beta}{\partial X}\delta X, \quad \delta\gamma=\frac{\partial\gamma}{\partial T}\delta T+\frac{\partial\gamma}{\partial X}\delta X;
\end{equation}
from which it is clear that the ratio $\frac{\delta\gamma}{\delta\beta}$ need not always be positive. If it is accepted that significant compositional variations exist in the lower mantle, then it is not immediately clear whether LLSVPs will be light or dense, and hence the role they play in mantle dynamics. To address this latter point what is needed are direct estimates of their density, and it free oscillations that have the potential to provide this information.

The first study of lateral density variations in the lower mantle using free oscillation data was undertaken in the late 90s by Miaki Ishii and Jeroen Tromp then both at Harvard. This study was based on a slight elaboration the first-order splitting theory discussed earlier in the lecture, and the resulting models showed that LLSVPs are denser than the surrounding mantle. This is consistent with the presence of compositional variations, and hence leads to interesting questions about what these variations are, and how they were produced. Do, for example, LLSVPs represent a "slab graveyard" comprised of compositionally distinct material accumulated at the base of the mantle through subduction. Or are they instead ancient iron-enriched features produced as the Earth differentiated into a mantle and core? Such questions continue to engage a wide range of Earth Scientists. But the density of LLSVPs is actually not a settled issue. This is because later work has shown that the simple first-order degenerate perturbation theory applied is not sufficiently accurate to model the effect of lateral density variations. Implicit within this theory is the assumption that distinct multiplets do not interact with each other on the addition of laterally varying structure. This is not true, with the resulting errors being especially significant for closely spaced multiplets. A better approach is to use normal mode coupling theory, this being based on the expansion of the eigenfunctions in a 3D earth model in the form
\begin{equation}
\mathbf{s}=\sum_{k'm'}\langle k'm'|P_{0}|\mathbf{s}\rangle|k'm'\rangle ,
\end{equation}
where $|km\rangle$ are eigenfunctions in a spherical reference model. Putting such an expansion into the rotating eigenvalue problem we arrive at
\begin{equation}
\sum_{k'm'}[-\omega^{2}\langle km|P|k'm'\rangle+i\omega\langle km|W|k'm'\rangle+\langle km|H|k'm'\rangle]\langle k'm'|P_{0}|\mathbf{s}\rangle=0 ,
\end{equation}
for all (k,m). This is an algebraic non-linear eigenvalue problem for the expansion coefficients $\langle k'm'|P_{0}|\mathbf{s}\rangle$ relative to the chosen basis. Such an equation cannot be solved directly because the matrices and vectors involved are infinite-dimensional. But it is possible to truncate the basis so as to obtain a finite-dimensional problem that can be solved numerically, and by taking sufficiently many terms in the expansion the results can be made as accurate as desired. The computational task here is fairly substantial, and it has only been in the past decade or so that such methods have become viable. It is through such calculations that the earlier first-order approaches have been shown to be inadequate, but there has yet to be an application of the newer theory to re-invert for lateral density variations and check if Ishii and Tromp's conclusions remain valid. A rather similar approach was, however, applied to observations of solid Earth tides in 2017 led by Harriet Lau\footnote{Then a PhD student and Harvard, and now a Professor at Brown University.}. As with free oscillations, solid Earth tides are sensitive to lateral density variations because of the large length-scales of the associated deformation. The result of Lau's study was that LLSVPs do indeed appear to be dense, but these conclusions remain somewhat tentative and work in this area is far from finished!

\section*{What you need to know and be able to do}
\begin{itemize}
    \item[(i)] How eigenfunctions of a spherically symmetric model can be classified in terms of mode type, degree, overtone number, and order. You do not need to learn the definition of vector spherical harmonics, and will not be asked to perform detailed calculations with them.
    \item[(ii)] How to solve the toroidal mode eigenvalue problem in a homogeneous earth model using spherical Bessel functions see the second problem set. Note that you might be asked to perform such calculations in the exam, but would be provided with the necessary information about special functions.
    \item[(iii)] How degenerate perturbation theory can be used to understand mode splitting due to rotation and lateral heterogeneity. You should know qualitatively what selections rules are and why they are useful, but you are not expected to remember their specific form or know how to derive them.
    \item[(iv)] What the diagonal sum rule says (but not the derivation), and how it was applied in determining spherically averaged density variations.
    \(v) The idea that convection in the mantle might be driven in part by compositional variations, and how this complicates the dynamic interpretation of tomographic velocity models.
\end{itemize}

\end{document}
